\begin{figure}[!htbp]
    \centering
    \begin{tikzpicture}[
        box/.style={draw, rounded corners=3pt, minimum height=0.7cm, align=center, font=\small},
        databox/.style={draw, rounded corners=3pt, minimum height=0.7cm, align=center, fill=black!8, font=\small},
        reprbox/.style={draw, rounded corners=3pt, minimum height=0.7cm, align=center, fill=black!5, font=\small},
        arr/.style={->, >=stealth, thick},
        lbl/.style={font=\footnotesize\bfseries},
        catlbl/.style={font=\small\bfseries, anchor=east},
    ]

    % 分类标签
    \node[catlbl] at (-0.3, 1.5) {连续表征};
    \node[catlbl] at (-0.3, -1.5) {离散表征};

    % 分隔线
    \draw[dashed, black!30] (-0.1, 0) -- (11.5, 0);

    % 梅尔频谱
    \node[reprbox, minimum width=2cm] (mel) at (1.6, 2.4) {梅尔频谱};
    \node[box, minimum width=1.6cm] (vocoder1) at (5.4, 2.4) {声码器};
    \node[databox, minimum width=1.6cm] (wav1) at (8.6, 2.4) {语音波形};
    \draw[arr] (mel) -- node[above, font=\scriptsize, text=black!60] {时频矩阵} (vocoder1);
    \draw[arr] (vocoder1) -- (wav1);

    % VAE
    \node[reprbox, minimum width=2cm] (vae) at (1.6, 0.8) {VAE潜在特征};
    \node[box, minimum width=1.6cm] (vaedec) at (5.4, 0.8) {VAE解码器};
    \node[databox, minimum width=1.6cm] (wav2) at (8.6, 0.8) {语音波形};
    \draw[arr] (vae) -- node[above, font=\scriptsize, text=black!60] {低维潜在向量} (vaedec);
    \draw[arr] (vaedec) -- (wav2);

    % 单码本
    \node[reprbox, minimum width=2cm] (single) at (1.6, -0.8) {单码本token};
    \node[box, minimum width=1.6cm] (synth1) at (5.4, -0.8) {语音合成器};
    \node[databox, minimum width=1.6cm] (wav3) at (8.6, -0.8) {语音波形};
    \draw[arr] (single) -- node[above, font=\scriptsize, text=black!60] {一维token序列} (synth1);
    \draw[arr] (synth1) -- (wav3);

    % 多码本
    \node[reprbox, minimum width=2cm] (multi) at (1.6, -2.4) {多码本token};
    \node[box, minimum width=1.6cm] (synth2) at (5.4, -2.4) {语音合成器};
    \node[databox, minimum width=1.6cm] (wav4) at (8.6, -2.4) {语音波形};
    \draw[arr] (multi) -- node[above, font=\scriptsize, text=black!60] {多层token矩阵} (synth2);
    \draw[arr] (synth2) -- (wav4);

    \end{tikzpicture}
    \bicaption{\enspace 用于生成的语音表征分类：连续表征（梅尔频谱和VAE潜在特征）与离散表征（单码本和多码本token）}{\enspace Taxonomy of speech representations for generation: continuous representations (mel spectrogram and VAE latent features) and discrete representations (single-codebook and multi-codebook tokens)}
    \label{fig:rw_repr_gen}
\end{figure}
